\documentclass[10pt]{article}
\usepackage[margin=2cm]{geometry}
\usepackage{amsmath,amssymb,booktabs}
\newcommand{\Tr}{\operatorname{Tr}}
\title{Batched (CPU/GPU) entanglement pipeline for qubit networks:\\ method, validation and comparison with QuTiP}
\author{}\date{}
\begin{document}\maketitle\vspace{-2em}

\section*{1. What was implemented}
For a graph $G$ on $n$ qubits, $H=\sum_{(i,j)\in E}h_{ij}$, every stage is vectorized over the whole time grid ($T=5001$) and written against an array module, so the same code runs on NumPy and on CuPy (GPU).
\begin{enumerate}\itemsep0pt
\item \textbf{Evolution.} One eigendecomposition $H=VEV^\dagger$, then $\psi(t)=V e^{-iEt}V^\dagger\psi_0$ for all $t$ at once (exact; no ODE solver).
\item \textbf{Partial trace.} Reshape $\psi(t)$ to a $d_S\times d_{\bar S}$ matrix $A_t$; $\rho_S(t)=A_tA_t^\dagger$.
\item \textbf{Concurrence.} $C=\max(0,\lambda_1-\lambda_2-\lambda_3-\lambda_4)$, $\lambda_i$ the square roots of the eigenvalues of $\sqrt\rho\,\tilde\rho\sqrt\rho$, $\tilde\rho=(\sigma_y\!\otimes\!\sigma_y)\rho^*(\sigma_y\!\otimes\!\sigma_y)$ (Hermitian form: only batched \texttt{eigh}).
\item \textbf{PPT test and negativity bound.} Per cut $M$: $\lambda_{\min}(\rho^{T_M})\ge0$ (PPT $\Rightarrow$ GMN$=0$); $N_M=\sum|\text{neg.\ eig}(\rho^{T_M})|$, $\mathrm{GMN}\le\min_M N_M$.
\item \textbf{Meyer--Wallach.} $Q=2\big(1-\tfrac1n\sum_i\Tr\rho_i^2\big)$.
\item \textbf{GMN (SDP)} for $k$-qubit subsets, $k=3,4,5$: Secs.~2--3. (Guhne--Seevinck witness, $C_{\rm GME}$ and the negativity bound have no QuTiP counterpart and are not compared.)
\end{enumerate}

\section*{2. GMN SDP, certificate and batched ADMM solver}
For $k$ qubits ($d=2^k$) let $\mathcal M$ be the $K=2^{k-1}-1$ bipartitions ($K=3,7,15$):
\begin{equation}
p^\star=\min_{W,P_M,Q_M}\operatorname{Re}\Tr(W\rho)\ \ \text{s.t.}\ \ W=P_M+Q_M^{T_M},\ 0\preceq P_M,Q_M\preceq I\ (M\in\mathcal M),\qquad \mathrm{GMN}=\max(0,-p^\star).
\end{equation}
\textbf{Dual certificate.} Since $\min_{0\preceq P\preceq I}\Tr(\Lambda P)=\nu(\Lambda)$ (sum of negative eigenvalues), any Hermitian $\{\Lambda_M\}$ with $\sum_M\Lambda_M=\rho$ gives the rigorous bound
$\mathrm{GMN}\le U(\Lambda)=\sum_M\big[|\nu(\Lambda_M)|+|\nu(\Lambda_M^{T_M})|\big]$.

\textbf{ADMM.} With $x=(W,\{P_M\},\{Q_M\})$, $Ax=(P_M+Q_M^{T_M}-W)_M$, $c=(\rho,0,\dots)$: as $T_M^2=\mathrm{id}$, $AA^\dagger=2I+\mathbf 1\mathbf 1^{\!\top}\!\otimes I$, so $(AA^\dagger)^{-1}=\tfrac12(I-\tfrac{J}{2+K})$ and $\Pi_{\ker A}(x)=x-A^\dagger(AA^\dagger)^{-1}Ax$ needs only elementwise operations and partial transposes. $W$ is unconstrained, hence $z_W=\hat x_W$ and its scaled dual is identically zero (not stored). With penalty $\varrho$, relaxation $\alpha$, scaled dual $u$:
\begin{align}
x^{+}&=\Pi_{\ker A}(z-u-c/\varrho),\qquad \hat x=\alpha x^{+}+(1-\alpha)z,\\
z^{+}&=\Pi_{\mathcal K}(\hat x+u)\ \ (\text{clip eigenvalues of }P_M,Q_M\text{ to }[0,1]),\qquad u^{+}=u+\hat x-z^{+}.
\end{align}
\textbf{Certified stopping.} The multiplier of the $x$-step, $\Lambda_M=\varrho y_M$ ($y=(AA^\dagger)^{-1}A(\cdot)$), satisfies $\sum_M\Lambda_M=\rho+\varrho(x_W-z_W)$, so after the Frobenius projection onto $\sum_M\Lambda_M=\rho$ it gives $U$. The primal estimate is $L=\max(0,-\operatorname{Re}\Tr W\rho)$. A state stops when $U-L\le10^{-7}$ and the primal residual $\le10^{-8}$; stopped states are removed from the batch. A batch of $B$ states costs $2KB$ batched $d\times d$ \texttt{eigh}s per iteration.
\textbf{Defaults:} $\varrho=0.003$, $\alpha=1.0$, 40000 iterations (tuned on hard states: Table~2).

\textbf{Decision at threshold $\tau=10^{-6}$:} \emph{zero} if $U\le\tau$ (rigorous); \emph{positive} if $L-(U-L)>\tau$ (margin rule, validated empirically, \emph{not} rigorous since $L$ is a primal estimate); otherwise \emph{near threshold}, reported with the bracket $[L,U]$.

\section*{3. Comparison with QuTiP}
ex$|{+}\rangle$ on the complete graph $K_n$, $T=5001$, Tesla V100, one CPU core; seconds, GPU inputs device-resident. All implementations agree with the QuTiP versions to $\le2\!\times\!10^{-8}$ (concurrence), $\le10^{-10}$ (other stages).
\begin{center}\small\textbf{Table 1} ($n=5$)\\[2pt]\begin{tabular}{lrrrr}\toprule
Stage & ours NumPy & ours CuPy & QuTiP alternative & QuTiP / CuPy\\\midrule
Evolution & 0.0099 & 0.0012 & 0.123 (\texttt{sesolve}); 34.0 (qutip-jax, GPU) & 100$\times$; 28000$\times$\\
Partial trace, all pairs & 0.035 & 0.0013 & 9.02 (\texttt{Qobj.ptrace} loop) & 7000$\times$\\
Concurrence, all pairs & 3.96 & 0.019 & 8.46 (\texttt{qt.concurrence} loop) & 440$\times$\\
PPT min-eig, all triples, 3 cuts & 1.27 & 0.052 & 11.6 (\texttt{partial\_transpose} loop) & 220$\times$\\
Meyer--Wallach $Q(t)$ & 0.017 & 0.0017 & 5.00 (\texttt{ptrace} loop) & 2900$\times$\\\bottomrule
\end{tabular}\end{center}
QuTiP times grow roughly linearly with $n$ (per-state Python loops); CuPy barely changes (the $n=3,4$ columns are in the V100 results file in \texttt{gpu\_bench}). \texttt{qutip-cuquantum} failed on the V100 (\texttt{CUTENSOR\_ERROR}); on a T4 (weak FP64) CuPy is slower than NumPy for the small stages, so the GPU type matters.

\textbf{GMN per solve} (cvxpy+SCS, one core, vs GPU ADMM; QuTiP has no GMN solver). Production uses $\sim\!40$ SCS workers ($\approx\!17$ ms per triple), so the GPU gain for triples is small; it is $\sim\!30\times$ for 4-qubit states, and the GPU is the only practical option for 5-qubit states. Times from an earlier fixed-iteration version.
\begin{center}\small\begin{tabular}{lccc}\toprule
 & $k=3$ ($8\!\times\!8$) & $k=4$ ($16\!\times\!16$) & $k=5$ ($32\!\times\!32$)\\\midrule
cvxpy+SCS, 1 core & 0.70 s & 160 s & infeasible\\
GPU ADMM & 0.053 s & 0.135 s & 1.1 s\\\bottomrule
\end{tabular}\end{center}

\section*{4. Validation and design choices}
\textbf{Tuning} (Table~2; K4, ex$|{+}\rangle$, converged = certified gap $\le10^{-7}$ within 40000 iterations): smaller $\varrho$ and $\alpha=1$ help the slowly converging states; $\varrho\ge0.1$ converges almost none.
\begin{center}\small\textbf{Table 2}\\[2pt]\begin{tabular}{lcc}\toprule
Set & old ($\varrho{=}0.03,\alpha{=}1.5$) & new ($\varrho{=}0.003,\alpha{=}1.0$)\\\midrule
3-qubit subset (400 states) & 95.5\% & 99.5\%\\
4-qubit whole state (150 states) & 67.3\% & 87.3\%\\\bottomrule
\end{tabular}\end{center}
\textbf{Against CLARABEL.} GHZ- and W-state sweeps through $\mathrm{GMN}=0$ ($k=3$: 46 states each; $k=4$: 6 of 46): no decision mismatches, max error $2.5\!\times\!10^{-8}$ (GHZ$_3$), $1.3\!\times\!10^{-8}$ (W$_3$), $9.6\!\times\!10^{-9}$ (GHZ$_4$). On the 60 real states with the smallest negativity bound per case ($\mathrm{GMN}\sim10^{-8}$--$10^{-5}$), CLARABEL never exceeded $U$ (6 settings, 45 comparisons), and no certified zero was wrong. The reference itself is unreliable there: every CLARABEL/SCS solve returned \texttt{optimal\_inaccurate}, and one input gave values from $0$ to $2.9\!\times\!10^{-7}$ across tolerances/solvers, so the $10^{-7}$--$10^{-6}$ range cannot be arbitrated by the CPU solvers (this may also affect the SCS-based phase-1 results at the margin; not checked). $L$ vs CLARABEL: below it for the old settings, above it by up to $8\!\times\!10^{-7}$ (3-qubit) for the new ones; for 4-qubit near-threshold states by up to $10^{-5}$ in one case that was not resolved, and none of these converged. Hence the margin rule in Sec.~2.

\textbf{Review of solver design} (external suggestions, tested): dropping $W_u$ is exact; the choice of multiplier estimate ($y$, $u_P$, averaged) makes no difference; the primal residual $\|x-z\|$ bounds the affine violation of $z$ within a factor $\le 3$; masking instead of compaction is 60--80\% slower; check intervals 25--500 change time by $<10\%$; per-problem adaptive $\varrho$ helps easy states by 3--10\% and does not help the stuck ones (so it is off by default).
\end{document}
